\documentclass[10pt,a4paper]{amsart}
\usepackage[margin=19mm]{geometry}
\usepackage{amsmath,amssymb,mathtools}
\usepackage[scr=rsfs,cal=euler]{mathalfa}
\usepackage{xfrac}
\usepackage[unicode,hidelinks,bookmarks=false]{hyperref}
\newcommand{\ie}{i.e.\ }
\newcommand{\cf}{cf.\ }
\newcommand{\resp}{respectively\ }
\newcommand{\Subsection}{\subsection}
\providecommand{\nobreakdash}{\nobreak\hbox{-}}
\setlength{\emergencystretch}{2em}
\multlinegap0pt
\usepackage{bm}
\usepackage{bbm}
\usepackage{dsfont}
\usepackage{xspace}
\usepackage{cancel}
\usepackage{mathtools}
\usepackage{amsthm}
\makeatletter
\@ifundefined{theorem}{\newtheorem{theorem}{Theorem}}{}
\makeatother
\newcommand{\df}{\dfrac}
\title[An Arithmetic Sum Associated with the Classical Theta Functi]{An Arithmetic Sum Associated with the Classical Theta Function}
\author{Bruce C. Berndt, Raghavendra N. Bhat, Jeffrey L. Meyer, Likun Xie, Alexandru Zaharescu}
\date{Source: arXiv:2501.03234v3 (CC BY 4.0)}
\begin{document}
\maketitle

\noindent\emph{Source and licence.} An Arithmetic Sum Associated with the Classical Theta Function. Version arXiv:2501.03234v3 (CC BY 4.0), distributed under the Creative Commons Attribution 4.0 International licence, as recorded in the arXiv licence metadata for this exact version and shown on the abstract page https://arxiv.org/abs/2501.03234v3. Full rights notice: licenses/arxiv-2501.03234-CC-BY-4.0.txt.\par\smallskip

\noindent\textit{Editorial scope.} Counted unit: the Theorem whose source label is \texttt{theorem2} in the article (source0.tex lines 867--875, source label \texttt{theorem2}), delivered together with the article's own complete original proof of it (source0.tex lines 877--937, one \texttt{proof} environment of 61 lines). No mathematics is written, repaired or re-derived: statement and proof are the article's own text line for line. Required context retained uncounted: S\_k\_expression (lines 695--699); r (lines 701--703); m\_jk\_expression (lines 721--724). Every definition, display and label of those lines is retained unchanged. Omitted from this package and not redistributed: 1--694, 700--700, 704--720, 725--866, 876--876, 938--1611. Nothing inside a retained excerpt is omitted or altered: every retained body is delivered exactly as the article's lines are written, so no omission receipt of any kind accompanies this package. Citations: the retained excerpts carry no citation command at all, so no bibliography entry of the article is transcribed and no reference is redistributed. Reference adaptation: none was needed and none was made; every reference inside the delivered unit resolves inside it, so no unresolved reference or citation is produced. Printed slips kept literal: no printed word, symbol, hypothesis or number of the counted statement or of its proof was altered. Layout and class: the article's own class is \texttt{amsart} with packages including latexsym, graphicx, float, fancyhdr, amssymb, mathtools, tabularx, comment, hyperref, amsmath, amsthm; the wrapper is the lane's pinned \texttt{amsart} editorial wrapper at 10pt on A4, which restates the theorem-like environments the retained text needs and adds the article's own macro definitions that the retained text needs (\texttt{\textbackslash df}), with extra wrapper packages (bm, bbm, dsfont, xspace, cancel, mathtools). The delivered numbering is the unit's own. Assets: no retained span contains a figure environment, a graphics-inclusion command, a PDF-page inclusion, a table or a TikZ picture; no image file is copied into this package, the package declares no asset dependency and no image file is redistributed with it. Licence: version arXiv:2501.03234v3 of this article is distributed under the Creative Commons Attribution 4.0 International licence, as recorded in the arXiv licence metadata for this exact version and shown on the abstract page https://arxiv.org/abs/2501.03234v3. A full rights notice is delivered at licenses/arxiv-2501.03234-CC-BY-4.0.txt. Characters: every retained excerpt is pure ASCII, so the delivered engine, XeLaTeX with its default Latin Modern text font, reports no missing character.
\begin{align}
	S(k)=\sum_{\substack{j=1\\j\text{ odd}}}^{k-1}\sum_{h=1}^{k-1} (-1)^{j+1+[hj/k]}
	=\sum_{\substack{j=1\\j\text{ odd}}}^{k-1}\sum_{h=1}^{k-1} (-1)^{[hj/k]}
	=\frac{(k-1)^2}{2}-2\sum_{\substack{j=1\\j\text{ odd}}}^{k-1}r(j,k), \label{S_k_expression}
\end{align}

\begin{equation}\label{r}
r(j,k)= \#\left\{h: 1\leq h\leq k-1;  \left[\df{hj}{k}\right] \text{ is odd }\right\}.
\end{equation}

\begin{align}
    r(j,k)= \sum_{h=1}^{k-1} \left(\left[\df{hj}{k}\right]-2\left[\df{hj}{2k}\right]\right)
    =\frac{(j-1)(k-1)}{2}-2\sum_{h=1}^{k-1}\left[\df{hj}{2k}\right].\label{m_jk_expression}
  \end{align}

\begin{theorem}\label{theorem2} If $k$ is an odd prime, then
 \begin{align*}
	S(k)
	&\geq -\frac{k-1}{2}+kH(k)-\frac{(k-1)(k+1)}{4},
\end{align*}
where, as $k\to\infty$,
 $$H(k)=\frac{1}{2}\log\left(2k\right)+\frac{\gamma}{2}+O\left(\frac{1}{k}\right),$$
where $\gamma$ denotes Euler's constant.
\end{theorem}

\begin{proof} Recall from \eqref{S_k_expression} the identity
\begin{align}\label{4.22}
	S(k)
	&=\frac{(k-1)^2}{2}-2\sum_{\substack{j=1\\j \text{ odd}}}^{k-1}r(j,k),
\end{align}
where  $r(j,k)$ is defined in \eqref{r}. From \eqref{m_jk_expression},
    \begin{align}\label{mmmm}
r(j,k)= \sum_{h=1}^{k-1} \left(\left[\frac{hj}{k}\right]-2\left[\frac{hj}{2k}\right]\right).
\end{align}
Thus, by \eqref{mmmm},
\begin{align}
 \lim_{k\to\infty}\frac{1}{k} r(j,k)&=\lim_{k\to\infty} \frac{1}{k}\sum_{h=1}^{k} \left(\left[\frac{hj}{k}\right]-2\left[\frac{hj}{2k}\right]\right)\nonumber\\
    &=\int _0^1\left(\left[ jx\right]-2\left[\frac{jx}{2}\right]\right)dx\nonumber\\
%&=\df{1}{j}\int _0^j\left(\left[ u\right]-2\left[\frac{u}{2}\right]\right)du\nonumber
    &=\df{1}{j}\left(\frac{j-1}{2}\right).\label{riemannsum}
\end{align}
To obtain the last equality in \eqref{riemannsum}, first note that
\begin{equation}\label{4.111}
\left[ jx\right]-2\left[\frac{jx}{2}\right]
=\begin{cases} 1,\quad &\text{ if } k+\frac{1}{2} \leq \frac{jx}{2} \leq k+1,  \\
0,&\text{ otherwise}.
\end{cases}
\end{equation}
 Next, subdivide the interval $[0,1] $ into subintervals of length $\frac{1}{j}$. From \eqref{4.111}, we see that the largest value of $k$ is $\frac{j-1}{2}$. The value of the integral is therefore $\frac{1}{j}\cdot\frac{j-1}{2}$, and consequently \eqref{riemannsum} follows.

From the limit in \eqref{riemannsum}, we are motivated to define an error term $\epsilon_j(k)$ by
\begin{equation}\label{4.12}
\epsilon_j(k):= \frac{1}{k} r(j,k) -\frac{1}{j}\left(\frac{j-1}{2}\right).
\end{equation}
 We derive an upper bound for $\epsilon_j(k).$  Define
 \begin{equation}\label{4.24}
  f(j,k):=[ jx]-2\left[\frac{jx}{2}\right],
  \end{equation}
  as in \eqref{4.111}.  Subdivide the interval $[0,1]$ into subintervals of length $\frac{1}{j}.$ Consider the sequence of fractions
 $$\frac{1}{k},\frac{2}{k},\dots,\frac{k-1}{k},$$
 and set
 $$a_n=\frac{n}{k},\quad 1\leq n \leq k-1.$$
 Consider a point
 $$\frac{n}{j}, \quad 1\leq n \leq j-1.$$
   Let $a_m$ be the closest fraction to $\frac{n}{j}$ to the left. Of course, $a_{m+1}$ is the closest fraction to $\frac{n}{j}$ on the right.  We note that $|a_{m+1}-a_m|=\frac{1}{k}$.  From \eqref{4.111} and \eqref{4.24}, we see that the contribution of each of these subintervals to $f(j,k)$ is $\frac12\cdot\frac{1}{k}$. Since there are $j-1$ of these contributions, we conclude that
 $$ \epsilon_j(k)\leq \dfrac{j-1}{2k}.$$
 Therefore, by \eqref{4.12},
 \begin{equation}\label{4.14}
 r(j,k)\leq \frac{k}{2}-\frac{k}{2j}+\frac{j-1}{2} .
\end{equation}

Finally, from \eqref{4.22} and \eqref{4.14}, we conclude that
\begin{align}\label{4.15}
	S(k)
	&\geq \frac{(k-1)^2}{2}-2\sum_{\substack{j=1\\j\text{ odd}}}^{k-1}\left(\frac{k}{2}-\frac{k}{2j}+\frac{j-1}{2}\right) \nonumber\\
 &=-\frac{k-1}{2}+k\sum_{\substack{j=1 \\j \text{ odd}}}^{k-1}\frac{1}{j}-\frac{(k-1)(k+1)}{4}.
\end{align}
Recall the familiar asymptotic formula for the harmonic sum, as $k\to\infty$,
$$ \sum_{j=1}^k\df{1}{j} = \log k +\gamma +O\left(\df{1}{k}\right), $$
where $\gamma$ denotes Euler's constant.  Hence, we deduce that, as $k\to\infty$,
\begin{align}\label{4.16}
\sum_{\substack{j=1\\j\text{ odd}}}^{k-1}\df{1}{j}= \sum_{j=1}^{k-1}\df{1}{j}- \df{1}{2}\sum_{j=1}^{(k-1)/2}\df{1}{j}=&\;\frac{1}{2}\log\left(2(k-1)\right)+\frac{\gamma}{2}+O\left(\frac{1}{k-1}\right)\notag\\
=&\;\frac{1}{2}\log(2k)+\frac{\gamma}{2}+O\left(\frac{1}{k}\right).
\end{align}
Together, \eqref{4.15} and \eqref{4.16} complete the proof of Theorem \ref{theorem2}.
\end{proof}

\end{document}
